\begin{houseviewbox}[结论：同业较低的远期 PE 说明市场已在定价高盈利，但不自动创造雅化上行]
截至 2026-07-23，可核验的三家锂资源/锂盐可比公司 2026E PE 为 11.71--13.51 倍，均值 12.79 倍、中位数 13.16 倍。雅化现价对应 House 基准 2026E PE 为 10.18 倍，表面上存在折价；但这一折价可以对应现金转换、成本与资源兑现的不确定性。因此可比倍数只校验估值区间，不能代替雅化自己的盈利、现金与分部证据。
\end{houseviewbox}

\section{相对 PE：统一价格日、不同业务质量}

可比组由天齐锂业、赣锋锂业和盛新锂能构成。它们与雅化同受锂周期影响，但资源禀赋、海外资产、产品结构、资本开支和会计口径不同；因此本报告没有把 12.79 倍直接乘以雅化 House EPS，也不以一个可比平均数给目标价。相对估值的合理作用是回答“基准 9.9 倍周期 PE 是否离谱”，而不是覆盖公司特有的现金和成本问题。

\begin{exhibitbox}[表 12-1：锂产业链可比公司 2026E PE]
\begin{tabularx}{\exhibitboxwidth}{L{3.0cm}L{2.0cm}R{2.3cm}R{2.3cm}X}
\toprule
公司 & 代码 & 2026E PE & 使用方式 & 可比性边界 \\
\midrule
天齐锂业 & 002466.SZ & 11.71x & 区间校验 & 资源、海外资产与权益结构不同 \\
赣锋锂业 & 002460.SZ & 13.51x & 区间校验 & 资源组合、下游布局不同 \\
盛新锂能 & 002240.SZ & 13.16x & 区间校验 & 产能与资源路径不同 \\
均值 / 中位数 & -- & 12.79x / 13.16x & 低权重市场标尺 & 不进入最终目标公式 \\
雅化现价 / House 基准 & 002497.SZ & 10.18x & 观察折价 & EPS、现金与 H2 均为 House 验证项 \\
\bottomrule
\end{tabularx}
\end{exhibitbox}

\section{从估值折价到资本纪律：市场给折价未必是错误}

雅化 2026Q1 货币资金加交易性金融资产为 26.03 亿元，短期借款、一年内到期非流动负债和长期借款合计 11.19 亿元，净现金代理约 14.84 亿元。以 17.75 元市值和该代理计算，EV 约 189.74 亿元；对应 House 熊/基/牛 EBITDA 24.5/33.5/40.3 亿元，EV/EBITDA 为 7.74/5.66/4.71 倍。这个结果并不高，但不能覆盖 2025A 和 2026Q1 负经营现金流所揭示的营运资本风险。

\begin{exhibitbox}[表 12-2：现价下的盈利和资本约束交叉检查]
\begin{tabularx}{\exhibitboxwidth}{L{3.0cm}R{2.0cm}R{2.0cm}R{2.0cm}X}
\toprule
指标 & 熊 & 基准 & 牛 & 解读 \\
\midrule
2026E EPS & 1.30 & 1.74 & 2.14 & 2026H1 预告加 H2 情景 \\
现价 PE & 13.64x & 10.18x & 8.28x & 先看利润是否可兑现 \\
2026E EBITDA & 24.5 & 33.5 & 40.3 & House 交叉检查输入 \\
现价 EV/EBITDA & 7.74x & 5.66x & 4.71x & 不作为目标价锚 \\
2026E OCF & -1.0 & 2.0 & 6.0 & 现金比倍数更具判别力 \\
自由现金流 & -4.0 & -1.5 & 2.0 & 未因净现金而提高目标价 \\
\bottomrule
\end{tabularx}
\end{exhibitbox}

\section{外部卖方与 House 的真正分歧}

三份完整外部报告的 2026E EPS 为 2.01、2.23、2.63 元，均高于 House 基准 1.74 元。东吴 2026-07-07 的 42 元目标使用 16 倍 2026E PE；国信与开源没有可用目标价。我们不把三份报告的预测平均为“共识”：报告日期、假设口径和目标价可得性不同，且目标价只有一个有效观测。它们的价值是明确市场愿意在什么盈利与倍数条件下给予更高估值，正是 House 需要通过 H1 与 H2 证明的上行路径。

\begin{sourcequalitybox}[相对估值的行动纪律]
若 H1 现金、单位经济和 H2 量价成本支持 House 基准，雅化对可比组的折价可被重新评估；若只看到利润预告而未看到现金和成本，低 PE 是风险补偿而非安全边际。故最终目标价仍由概率加权内在价值主导，市场和外部锚合计不改变基本面判断。
\end{sourcequalitybox}

\sourcenote{可比公司 2026E PE 快照；公司 2026Q1 报告；国信、开源、东吴原始报告及外部预测差异表。}
